\par\Needspace{7\baselineskip}\bigskip
\begin{center}
\pdfbookmark[2]{Olbers to Gauss, Bremen, 7 March 1816}{pallas-brief-18160307}
\textsc{Olbers} to \textsc{Gauss}, Bremen, 7 March 1816.
\end{center}
\noindent\hspace*{\parindent}\SourceLeaf{433}\{. . . . .\quad For now I hasten only to give you further details of the Paris prize problem, in case you have not yet heard. The decision on the prize for the perturbations of the planets, especially Pallas, has again been postponed until 1817. The papers must, however, be submitted before 1 October 1816.
\par The prize question was: »The theory of planets whose eccentricity and inclination are too great for us to be able to calculate their perturbations accurately by the methods already known.« The Class asks for no numerical application, only analytical formulas, but arranged in such a way that a skilled calculator can apply them with confidence either to the planet Pallas or to another of the planets newly discovered or still to be discovered. — Only two papers were received, and their authors paid insufficient attention to the aim the Class had stated in the prize announcement. Both (especially one of them) still passed over too many analytical developments, which mathematicians would first have to carry out to be able to understand and assess the solution of the problem they have given. They have been too remiss in coming down to the level of the calculator who might now wish to construct tables for Pallas or some other planet. The supplements they submitted at various times are far from removing all these difficulties. From these supplements and the notes submitted by the anonymous authors, the Class saw that they had not had time to undertake all the necessary developments. It also considers that the same cause may have kept other mathematicians with the ability and skill to treat this difficult subject from entering for the prize, and has therefore postponed the award until January 1817. — The prize is double, a gold medal worth 6000 francs.\quad\mbox{. . . . .\}}
\par
